\chapter{Certified spin chains and ladders}
\label{chap:models}

The reflection-positive ring criterion (Theorem~\ref{thm:rp-gap}) applies across spin-1 and
spin-\(\tfrac12\) chains and ladders, both for Hamiltonians that are reflection positive in a
natural real basis (the computational \(Z\)-basis for the Ising models, the Cartesian basis of
Definition~\ref{def:spin-one-generators} for spin~1) and for antiferromagnetic models that become
reflection positive after a bipartite sublattice rotation (Section~\ref{sec:sublattice-rot} and
Theorems~\ref{thm:conj-hamiltonian}--\ref{thm:conj-spectrum}). Each model is formulated as a
reflection-positive ring (Definition~\ref{def:rpring}) together with its certified parameter
points, seed inverse temperature \(\beta_0\), and rational doubling certificate
(Definition~\ref{def:chain-certificate}), yielding a uniform gap bound \(\log r / \beta_0\) at
every point by Theorem~\ref{thm:rp-gap}. The only open hypothesis for each model is a single lemma
stating its two finite-size seed inequalities~\eqref{eq:rp-seed-hyp} at every certified parameter
point, which are verified numerically from complete eigenvalue enclosures of the seed
rings. Ring lengths count unit cells.

\section{Local spin operators, sublattice rotations, and cell reflection}
\label{sec:sublattice-rot}

The spin-\(\tfrac12\) and bond-alternating (dimerized) models share three local operators: the real
spin-\(\tfrac12\) generators \(P_\alpha\) (Definition~\ref{def:local-spin-half}), the bipartite
sublattice rotation \(e^{i\pi S^y}\), which on one site of a bond turns the antiferromagnetic
exchange interaction into minus a sum of tensor squares of real operators~\eqref{eq:spin-half-rot}
(Lemma~\ref{lem:spin-half-rotation}), and the two-site cell swap \(F\) used as the internal
reflection on dimerized rings (Definition~\ref{def:cell-swap}).

\begin{definition}[Pauli matrices and real spin-\(\tfrac12\) generators]
  \label{def:local-spin-half}
  \lean{CertifiedGaps.pauliZ, CertifiedGaps.pauliX, CertifiedGaps.pauliZ_transpose,
    CertifiedGaps.pauliX_transpose, CertifiedGaps.SpinHalf.P, CertifiedGaps.SpinHalf.P_symm}
  \leanok
  Let \(Z = \operatorname{diag}(1, -1)\) and \(X\) be the real symmetric Pauli matrices, and let
  \(P_\alpha = (S^x, iS^y, S^z)\) be the real spin-\(\tfrac12\) generators, each of which is
  symmetric or antisymmetric.
\end{definition}

\begin{lemma}[Bipartite sublattice rotation of the spin-\(\tfrac12\) exchange interaction]
  \label{lem:spin-half-rotation}
  \lean{CertifiedGaps.SpinHalf.spinHalf, CertifiedGaps.SpinHalf.rot,
    CertifiedGaps.SpinHalf.rot_mul_conjTranspose, CertifiedGaps.SpinHalf.rot_conj_spinHalf,
    CertifiedGaps.SpinHalf.Pc_eq}
  \leanok
  \uses{def:local-spin-half}
  The \(\pi\)-rotation \(R = e^{i\pi S^y} = \bigl(\begin{smallmatrix} 0 & 1 \\ -1 & 0
  \end{smallmatrix}\bigr)\)
  is unitary and satisfies \(R S^\alpha R^\dagger = s_\alpha S^\alpha\) with \(s = (-1, 1, -1)\).
  Since the real generators of Definition~\ref{def:local-spin-half} satisfy
  \(P_\alpha = c_\alpha S^\alpha\) with \(c = (1, i, 1)\), rotating one site of the exchange
  interaction \(\mathbf S \cdot \mathbf S\) transforms it into
  \begin{equation}
    \label{eq:spin-half-rot}
    (1 \otimes R)\,(\mathbf S \cdot \mathbf S)\,(1 \otimes R^\dagger)
    = -\sum_\alpha P_\alpha \otimes P_\alpha.
  \end{equation}
\end{lemma}
\begin{proof}
  \leanok
  Evaluate both sides of~\eqref{eq:spin-half-rot} entrywise on \(2 \times 2\) matrices.
\end{proof}

\begin{definition}[Two-site cell reflection operator]
  \label{def:cell-swap}
  \lean{CertifiedGaps.cellSwap, CertifiedGaps.cellSwap_mul_kronecker_mul_cellSwap,
    CertifiedGaps.cellSwap_mul_cellSwap, CertifiedGaps.cellSwap_transpose,
    CertifiedGaps.cellSwap_conj_smul_sum_kronecker_self}
  \leanok
  The swap operator \(F\) exchanging the two sites of a unit cell is a real symmetric involution
  satisfying \(F (A \otimes B) F = B \otimes A\) and commuting with every intra-cell exchange term
  \(c \sum_k A_k \otimes A_k\). It serves as the internal cell reflection \(F\) of
  Definition~\ref{def:rpring} for the dimerized models.
\end{definition}

\section{Spin-1 chain with single-ion anisotropy}

Adding a positive single-ion anisotropy \(D \sum_j (S_j^z)^2\) to the spin-1 Heisenberg chain
\eqref{eq:heisenberg-ham} favors the \(S^z = 0\) projection on each site. Beyond a quantum phase
transition near \(D \approx 1\), the chain leaves the Haldane phase of \(D = 0\) and enters the
gapped large-\(D\) phase, which contains both certified parameter points of
Definition~\ref{def:large-d-points}. Because the anisotropy is a real symmetric on-site operator,
reflection positivity in the Cartesian basis (Definition~\ref{def:spin-one-large-d}) is preserved,
yielding Theorem~\ref{thm:large-d-gap}.

\begin{definition}[Certified couplings for the single-ion anisotropic spin-1 chain]
  \label{def:large-d-points}
  \lean{CertifiedGaps.SpinOneLargeD.Point, CertifiedGaps.SpinOneLargeD.points}
  \leanok
  \uses{def:spin-one-large-d, def:chain-certificate}
  The anisotropic spin-1 chain \(\sum_j \mathbf S_j \cdot \mathbf S_{j+1} + D \sum_j (S^z_j)^2\) of
  Definition~\ref{def:spin-one-large-d} is certified at the two couplings \(D = 4\)
  (\(\beta_0 = 9/4\), \(r = 4.755\)) and \(D = 1603/500\) (\(\beta_0 = 381/125\),
  \(r = 1.0065429\)), each with seed length \(n_0 = 6\) sites.
\end{definition}

\begin{lemma}[Admissibility of the large-\(D\) doubling certificates]
  \label{lem:large-d-valid}
  \lean{CertifiedGaps.SpinOneLargeD.points_valid, CertifiedGaps.SpinOneLargeD.points_β₀_pos}
  \leanok
  \uses{def:large-d-points}
  Both rational doubling certificates of Definition~\ref{def:large-d-points} are admissible in the
  sense of Definition~\ref{def:chain-certificate}, and both satisfy \(\beta_0 > 0\).
\end{lemma}
\begin{proof}
  \leanok
  Verified by exact rational arithmetic from Definition~\ref{def:chain-certificate}.
\end{proof}

\begin{lemma}[Finite-size seed purities for the large-\(D\) chain]
  \label{lem:large-d-seeds}
  \lean{CertifiedGaps.SpinOneLargeD.seeds}
  \leanok
  \uses{def:large-d-points, def:rp-purities}
  At each coupling in Definition~\ref{def:large-d-points}, \(\Ps(6, \beta_0) \ge 1 - s_0\) and
  \(\Pt(12, \beta_0) \ge 1 - t_0\).
\end{lemma}
\begin{proof}
  These finite-size bounds~\eqref{eq:rp-seed-hyp} are certified from complete eigenvalue enclosures
  of the seed rings and enter the formalization as an explicit numerical hypothesis.
\end{proof}

\begin{theorem}[Uniform spectral gap in the large-\(D\) phase]
  \label{thm:large-d-gap}
  \lean{CertifiedGaps.SpinOneLargeD.gap}
  \leanok
  \uses{lem:large-d-valid, lem:large-d-seeds, thm:rp-gap}
  At each certified coupling of Definition~\ref{def:large-d-points}, every even ring of \(L \ge 6\)
  sites has a unique ground state and a spectral gap greater than \(\log r / \beta_0\), giving
  \(\Delta_L > 0.693\) at \(D = 4\) and \(\Delta_L > 0.00214\) at \(D = 1603/500\).
\end{theorem}
\begin{proof}
  \leanok
  Apply Theorem~\ref{thm:rp-gap} with Lemmas~\ref{lem:large-d-valid} and~\ref{lem:large-d-seeds}.
\end{proof}

\section{Bond-alternating (dimerized) spin-1 chain}

Alternating the exchange coupling between weak and strong bonds along the spin-1 chain
\eqref{eq:spin-one-dimerized-ham} breaks single-site translation invariance down to a two-site unit
cell. Grouping each weak bond into a two-site unit cell with the cell swap \(F\)
(Definition~\ref{def:cell-swap}) as its internal reflection makes the strong inter-cell bonds
reflection positive in the Cartesian basis (Definition~\ref{def:spin-one-dimerized}), yielding
Theorem~\ref{thm:spin-one-dimerized-gap}.

\begin{definition}[Dimerized spin-1 ring in the Cartesian basis]
  \label{def:spin-one-dimerized}
  \lean{CertifiedGaps.SpinOneDimerized.model, CertifiedGaps.SpinOneDimerized.Point,
    CertifiedGaps.SpinOneDimerized.points}
  \leanok
  \uses{def:rpring, def:spin-one-generators, def:cell-swap, def:chain-certificate}
  The bond-alternating spin-1 chain with Hamiltonian
  \begin{equation}
    \label{eq:spin-one-dimerized-ham}
    H_L = \sum_{j=0}^{L-1} (1-\delta)\, \mathbf S_{2j} \cdot \mathbf S_{2j+1} + (1+\delta)\,
    \mathbf S_{2j+1} \cdot \mathbf S_{2j+2}
  \end{equation}
  is formulated as a reflection-positive ring (Definition~\ref{def:rpring}) of weak-bond unit cells
  with bond operators \(C_\alpha = \sqrt{1+\delta}\,(1 \otimes G_\alpha)\) from
  Definition~\ref{def:spin-one-generators}, internal reflection \(F\) given by the cell swap of
  Definition~\ref{def:cell-swap}, and on-site term
  \(o = -(1-\delta) \sum_\alpha G_\alpha \otimes G_\alpha\). It is certified at dimerization
  \(\delta = 9/10\) with seed length \(n_0 = 2\) cells, \(\beta_0 = 7/2\), and \(r = 6.339\).
\end{definition}

\begin{lemma}[Admissibility of the dimerized spin-1 certificate]
  \label{lem:spin-one-dimerized-valid}
  \lean{CertifiedGaps.SpinOneDimerized.points_valid, CertifiedGaps.SpinOneDimerized.points_β₀_pos}
  \leanok
  \uses{def:spin-one-dimerized}
  The rational doubling certificate of Definition~\ref{def:spin-one-dimerized} is admissible and
  \(\beta_0 > 0\).
\end{lemma}
\begin{proof}
  \leanok
  Verified by exact rational arithmetic from Definition~\ref{def:chain-certificate}.
\end{proof}

\begin{lemma}[Finite-size seed purities for the dimerized spin-1 chain]
  \label{lem:spin-one-dimerized-seeds}
  \lean{CertifiedGaps.SpinOneDimerized.seeds}
  \leanok
  \uses{def:spin-one-dimerized, def:rp-purities}
  For the ring of Definition~\ref{def:spin-one-dimerized}, the seed purities satisfy
  \(\Ps(2, 7/2) \ge 1 - 0.04972\) and \(\Pt(4, 7/2) \ge 1 - 0.03206\).
\end{lemma}
\begin{proof}
  These finite-size bounds~\eqref{eq:rp-seed-hyp} are certified from complete eigenvalue enclosures
  of the seed rings and enter the formalization as an explicit numerical hypothesis.
\end{proof}

\begin{theorem}[Uniform spectral gap of the dimerized spin-1 ring]
  \label{thm:spin-one-dimerized-gap}
  \lean{CertifiedGaps.SpinOneDimerized.gap}
  \leanok
  \uses{lem:spin-one-dimerized-valid, lem:spin-one-dimerized-seeds, thm:rp-gap}
  Every even ring of \(L \ge 2\) unit cells of the dimerized spin-1 chain
  \eqref{eq:spin-one-dimerized-ham} at \(\delta = 9/10\) has a unique ground state and
  \(\Delta_L > \log 6.339 / (7/2) \approx 0.528\).
\end{theorem}
\begin{proof}
  \leanok
  Apply Theorem~\ref{thm:rp-gap} with Lemmas~\ref{lem:spin-one-dimerized-valid}
  and~\ref{lem:spin-one-dimerized-seeds}.
\end{proof}

\section{Two-leg transverse-field Ising ladder}

Coupling two ferromagnetic transverse-field Ising chains by ferromagnetic rungs produces a two-leg
ladder~\eqref{eq:two-leg-ising-ham} whose two-site rung cells are directly reflection positive in
the computational \(Z\)-basis (Definition~\ref{def:two-leg-ising} and
Theorem~\ref{thm:two-leg-ising-gap}).

\begin{definition}[Two-leg transverse-field Ising ladder Hamiltonian]
  \label{def:two-leg-ising}
  \lean{CertifiedGaps.TwoLegIsing.model, CertifiedGaps.TwoLegIsing.Point,
    CertifiedGaps.TwoLegIsing.points}
  \leanok
  \uses{def:rpring, def:local-spin-half, def:chain-certificate}
  The two-leg Ising ladder with Hamiltonian
  \begin{equation}
    \label{eq:two-leg-ising-ham}
    H_L = -\sum_{\text{legs}} Z_i Z_j - \sum_{\text{rungs}} Z_i Z_j - g \sum_i X_i
  \end{equation}
  is formulated as a reflection-positive ring (Definition~\ref{def:rpring}) of two-site rungs with
  bond operators \(Z \otimes 1\) and \(1 \otimes Z\) and on-site term
  \(o = -Z \otimes Z - g (X \otimes 1 + 1 \otimes X)\). It is certified at transverse field
  \(g = 3\) with seed length \(n_0 = 4\) rungs, \(\beta_0 = 3/2\), and \(r = 2.585\).
\end{definition}

\begin{lemma}[Admissibility of the two-leg Ising ladder certificate]
  \label{lem:two-leg-ising-valid}
  \lean{CertifiedGaps.TwoLegIsing.points_valid, CertifiedGaps.TwoLegIsing.points_β₀_pos}
  \leanok
  \uses{def:two-leg-ising}
  The rational doubling certificate of Definition~\ref{def:two-leg-ising} is admissible and
  \(\beta_0 > 0\).
\end{lemma}
\begin{proof}
  \leanok
  Verified by exact rational arithmetic from Definition~\ref{def:chain-certificate}.
\end{proof}

\begin{lemma}[Finite-size seed purities for the two-leg Ising ladder]
  \label{lem:two-leg-ising-seeds}
  \lean{CertifiedGaps.TwoLegIsing.seeds}
  \leanok
  \uses{def:two-leg-ising, def:rp-purities}
  For the ladder of Definition~\ref{def:two-leg-ising}, the seed purities satisfy
  \(\Ps(4, 3/2) \ge 1 - 0.1125\) and \(\Pt(8, 3/2) \ge 1 - 0.05162\).
\end{lemma}
\begin{proof}
  These finite-size bounds~\eqref{eq:rp-seed-hyp} are certified from complete eigenvalue enclosures
  of the seed rings and enter the formalization as an explicit numerical hypothesis.
\end{proof}

\begin{theorem}[Uniform spectral gap of the two-leg Ising ladder]
  \label{thm:two-leg-ising-gap}
  \lean{CertifiedGaps.TwoLegIsing.gap}
  \leanok
  \uses{lem:two-leg-ising-valid, lem:two-leg-ising-seeds, thm:rp-gap}
  Every even ring of \(L \ge 4\) rungs of the two-leg Ising ladder~\eqref{eq:two-leg-ising-ham} at
  \(g = 3\) has a unique ground state and \(\Delta_L > \log 2.585 / (3/2) \approx 0.633\).
\end{theorem}
\begin{proof}
  \leanok
  Apply Theorem~\ref{thm:rp-gap} with Lemmas~\ref{lem:two-leg-ising-valid}
  and~\ref{lem:two-leg-ising-seeds}.
\end{proof}

\section{Mixed-field Ising chain}

Adding a longitudinal field to the ferromagnetic transverse-field Ising chain breaks the global
\(\mathbb{Z}_2\) spin-flip symmetry, lifting the ground-state degeneracy of the ordered phase and
making the chain~\eqref{eq:mixed-ising-ham} non-integrable while keeping the ferromagnetic bond
interaction reflection positive in the computational \(Z\)-basis
(Definition~\ref{def:ising-longitudinal} and Theorem~\ref{thm:ising-longitudinal-gap}).

\begin{definition}[Mixed-field Ising ring Hamiltonian]
  \label{def:ising-longitudinal}
  \lean{CertifiedGaps.IsingLongitudinal.model, CertifiedGaps.IsingLongitudinal.Point,
    CertifiedGaps.IsingLongitudinal.points}
  \leanok
  \uses{def:rpring, def:local-spin-half, def:chain-certificate}
  The mixed-field Ising chain has Hamiltonian
  \begin{equation}
    \label{eq:mixed-ising-ham}
    H_L = -\sum_{j=0}^{L-1} Z_j Z_{j+1} - g \sum_{j=0}^{L-1} X_j - h \sum_{j=0}^{L-1} Z_j,
  \end{equation}
  represented as a reflection-positive ring (Definition~\ref{def:rpring}) with bond operator \(Z\)
  and on-site term \(-gX - hZ\), which is non-integrable for \(h \ne 0\). The formalization
  includes thirty certified grid points spanning \(g \in [1/4, 2]\) and \(h \in [1/20, 1]\),
  covering both the critical transverse field \(g = 1\) and the ferromagnetic regime \(g < 1\),
  each with seed length \(n_0 = 6\).
\end{definition}

\begin{lemma}[Admissibility of the mixed-field Ising certificates]
  \label{lem:ising-longitudinal-valid}
  \lean{CertifiedGaps.IsingLongitudinal.points_valid,
    CertifiedGaps.IsingLongitudinal.points_β₀_pos}
  \leanok
  \uses{def:ising-longitudinal}
  All thirty rational doubling certificates of Definition~\ref{def:ising-longitudinal} are
  admissible and every point satisfies \(\beta_0 > 0\).
\end{lemma}
\begin{proof}
  \leanok
  Verified by exact rational arithmetic from Definition~\ref{def:chain-certificate}.
\end{proof}

\begin{lemma}[Finite-size seed purities for the mixed-field Ising chain]
  \label{lem:ising-longitudinal-seeds}
  \lean{CertifiedGaps.IsingLongitudinal.seeds}
  \leanok
  \uses{def:ising-longitudinal, def:rp-purities}
  At each grid point in Definition~\ref{def:ising-longitudinal}, \(\Ps(6, \beta_0) \ge 1 - s_0\) and
  \(\Pt(12, \beta_0) \ge 1 - t_0\).
\end{lemma}
\begin{proof}
  These finite-size bounds~\eqref{eq:rp-seed-hyp} are certified from complete eigenvalue enclosures
  of the seed rings. The full thirty-point lemma enters as an explicit numerical hypothesis, while
  the seed at the first grid point \((g, h) = (1/4, 1/20)\) is proved unconditionally by exact
  integer arithmetic in Theorem~\ref{thm:ising-seeds}.
\end{proof}

\begin{theorem}[Uniform spectral gap of the mixed-field Ising ring]
  \label{thm:ising-longitudinal-gap}
  \lean{CertifiedGaps.IsingLongitudinal.gap}
  \leanok
  \uses{lem:ising-longitudinal-valid, lem:ising-longitudinal-seeds, thm:rp-gap}
  At each certified point of Definition~\ref{def:ising-longitudinal}, every even ring of
  \(L \ge 6\) sites of the mixed-field Ising chain~\eqref{eq:mixed-ising-ham} has a unique ground
  state and a spectral gap greater than \(\log r / \beta_0\), with certified gap lower bounds
  ranging from \(0.075\) at \((g, h) = (1, 1/20)\) to \(3.70\).
\end{theorem}
\begin{proof}
  \leanok
  Apply Theorem~\ref{thm:rp-gap} with Lemmas~\ref{lem:ising-longitudinal-valid}
  and~\ref{lem:ising-longitudinal-seeds}.
\end{proof}

\section{Spin-1/2 Heisenberg chain in a staggered field}

An antiferromagnetic spin-\(\tfrac12\) Heisenberg chain in a staggered magnetic field
\eqref{eq:staggered-ham} is not reflection positive in the standard \(S^z\) basis because its
exchange coupling has the antiferromagnetic sign. Applying the sublattice \(\pi\)-rotation
\(e^{i\pi S^y}\) on every odd site transforms each nearest-neighbor exchange
\(\mathbf S \cdot \mathbf S\) into \(-S^xS^x + S^yS^y - S^zS^z = -\sum_\alpha P_\alpha \otimes
P_\alpha\) by~\eqref{eq:spin-half-rot} and maps the staggered field to a uniform longitudinal
field, producing a reflection-positive ring (Definition~\ref{def:staggered}) with an identical
eigenvalue spectrum (Lemma~\ref{lem:staggered-physical-eq} and
Theorem~\ref{thm:staggered-physical-gap}).

\begin{definition}[Staggered-field chain in the reflection-positive frame]
  \label{def:staggered}
  \lean{CertifiedGaps.StaggeredHeisenberg.model, CertifiedGaps.StaggeredHeisenberg.Point,
    CertifiedGaps.StaggeredHeisenberg.points}
  \leanok
  \uses{def:rpring, def:local-spin-half, def:chain-certificate}
  In the sublattice-rotated frame, the staggered-field chain is the reflection-positive ring
  (Definition~\ref{def:rpring}) with bond operators \(P_\alpha\) from
  Definition~\ref{def:local-spin-half} and on-site term \(-h S^z\). It is certified at the five
  field strengths \(h = 1\), \(3/4\), \(1/2\), \(7/20\), and \(1/4\) on seed rings of \(n_0 = 8\)
  sites at inverse temperatures \(\beta_0 = 21/4\), \(21/4\), \(11/2\), \(23/4\), and \(25/4\).
\end{definition}

\begin{lemma}[Admissibility of the staggered-field certificates]
  \label{lem:staggered-valid}
  \lean{CertifiedGaps.StaggeredHeisenberg.points_valid,
    CertifiedGaps.StaggeredHeisenberg.points_β₀_pos}
  \leanok
  \uses{def:staggered}
  All five rational doubling certificates of Definition~\ref{def:staggered} are admissible and every
  point satisfies \(\beta_0 > 0\).
\end{lemma}
\begin{proof}
  \leanok
  Verified by exact rational arithmetic from Definition~\ref{def:chain-certificate}.
\end{proof}

\begin{lemma}[Finite-size seed purities for the staggered-field chain]
  \label{lem:staggered-seeds}
  \lean{CertifiedGaps.StaggeredHeisenberg.seeds}
  \leanok
  \uses{def:staggered, def:rp-purities}
  At each point in Definition~\ref{def:staggered}, \(\Ps(8, \beta_0) \ge 1 - s_0\) and
  \(\Pt(16, \beta_0) \ge 1 - t_0\).
\end{lemma}
\begin{proof}
  These finite-size bounds~\eqref{eq:rp-seed-hyp} are certified from complete eigenvalue enclosures
  of the seed rings and enter the formalization as an explicit numerical hypothesis.
\end{proof}

\begin{theorem}[Spectral gap of the staggered-field chain in the rotated frame]
  \label{thm:staggered-gap}
  \lean{CertifiedGaps.StaggeredHeisenberg.gap}
  \leanok
  \uses{lem:staggered-valid, lem:staggered-seeds, thm:rp-gap}
  At each certified field strength of Definition~\ref{def:staggered}, every even ring of
  \(L \ge 8\) sites in the reflection-positive frame has a unique ground state and a spectral gap
  greater than \(\log r / \beta_0\).
\end{theorem}
\begin{proof}
  \leanok
  Apply Theorem~\ref{thm:rp-gap} with Lemmas~\ref{lem:staggered-valid} and~\ref{lem:staggered-seeds}.
\end{proof}

\begin{definition}[Staggered-field Heisenberg Hamiltonian in the laboratory frame]
  \label{def:staggered-physical}
  \lean{CertifiedGaps.StaggeredHeisenberg.physicalHamiltonian,
    CertifiedGaps.StaggeredHeisenberg.physicalSpectrum, CertifiedGaps.StaggeredHeisenberg.rotation}
  \leanok
  \uses{def:ring-operators, lem:spin-half-rotation}
  In the laboratory frame, the staggered-field Heisenberg chain has Hamiltonian
  \begin{equation}
    \label{eq:staggered-ham}
    H_L = \sum_{j=0}^{L-1} \mathbf S_j \cdot \mathbf S_{j+1} - h \sum_{j=0}^{L-1} (-1)^j S^z_j
  \end{equation}
  and is related to the reflection-positive frame of Definition~\ref{def:staggered} by applying the
  sublattice rotation \(e^{i\pi S^y}\) on every odd site.
\end{definition}

\begin{lemma}[Sublattice unitary equivalence for the staggered-field chain]
  \label{lem:staggered-physical-eq}
  \lean{CertifiedGaps.StaggeredHeisenberg.physicalHamiltonian_eq,
    CertifiedGaps.StaggeredHeisenberg.physicalSpectrum_eq}
  \leanok
  \uses{def:staggered-physical, def:staggered, thm:conj-hamiltonian, thm:conj-spectrum}
  On an even ring, the laboratory-frame Hamiltonian~\eqref{eq:staggered-ham} is unitarily conjugate
  to the reflection-positive Hamiltonian of Definition~\ref{def:staggered} under the sublattice
  rotation, so their eigenvalue spectra agree.
\end{lemma}
\begin{proof}
  \leanok
  On an even ring every bond connects an odd site and an even site, so the sublattice rotation maps
  \(\mathbf S \cdot \mathbf S\) to \(-\sum_\alpha P_\alpha \otimes P_\alpha\)
  by~\eqref{eq:spin-half-rot} in Lemma~\ref{lem:spin-half-rotation} and turns the staggered field
  into the uniform longitudinal field. Theorems~\ref{thm:conj-hamiltonian}
  and~\ref{thm:conj-spectrum} then show via~\eqref{eq:conj-hamiltonian} that the two Hamiltonians
  are unitarily conjugate and share the same eigenvalue family.
\end{proof}

\begin{theorem}[Uniform spectral gap of the staggered-field Heisenberg ring]
  \label{thm:staggered-physical-gap}
  \lean{CertifiedGaps.StaggeredHeisenberg.physical_gap}
  \leanok
  \uses{lem:staggered-physical-eq, thm:staggered-gap}
  At each certified field strength, every even ring of \(L \ge 8\) sites of the staggered-field
  Heisenberg chain~\eqref{eq:staggered-ham} has a unique ground state and spectral gap
  \(> 1.1475, 0.8899, 0.5940, 0.3828, 0.2139\) at \(h = 1, 3/4, 1/2, 7/20, 1/4\).
\end{theorem}
\begin{proof}
  \leanok
  Combine Lemma~\ref{lem:staggered-physical-eq} with Theorem~\ref{thm:staggered-gap}.
\end{proof}

\section{Dimerized spin-1/2 Heisenberg chain}

Grouping a bond-alternating spin-\(\tfrac12\) Heisenberg chain~\eqref{eq:dimerized-half-ham} into
two-site unit cells across the strong bonds and rotating the second spin of each cell makes the
inter-cell bonds reflection positive, with the cell swap \(F\) (Definition~\ref{def:cell-swap}) as
the internal reflection, and turns the intra-cell exchange into a real symmetric on-site term that
commutes with \(F\) (Definition~\ref{def:dimerized}, Lemma~\ref{lem:dimerized-physical-eq}, and
Theorem~\ref{thm:dimerized-physical-gap}).

\begin{definition}[Dimerized spin-\(\tfrac12\) ring in the reflection-positive frame]
  \label{def:dimerized}
  \lean{CertifiedGaps.DimerizedHeisenberg.model, CertifiedGaps.DimerizedHeisenberg.Point,
    CertifiedGaps.DimerizedHeisenberg.points}
  \leanok
  \uses{def:rpring, def:local-spin-half, def:cell-swap, def:chain-certificate}
  In the sublattice-rotated frame, the dimerized chain is a reflection-positive ring
  (Definition~\ref{def:rpring}) of strong-bond unit cells with bond operators
  \(C_\alpha = \sqrt{1-\delta}\,(1 \otimes P_\alpha)\), internal reflection \(F\) given by the cell
  swap of Definition~\ref{def:cell-swap}, and on-site term
  \(o = -(1+\delta) \sum_\alpha P_\alpha \otimes P_\alpha\), certified at the three dimerizations
  \(\delta = 1/2, 3/10, 1/4\) with seed length \(n_0 = 4\) cells and inverse temperatures
  \(\beta_0 = 6, 13/2, 29/4\).
\end{definition}

\begin{lemma}[Admissibility of the dimerized spin-\(\tfrac12\) certificates]
  \label{lem:dimerized-valid}
  \lean{CertifiedGaps.DimerizedHeisenberg.points_valid,
    CertifiedGaps.DimerizedHeisenberg.points_β₀_pos,
    CertifiedGaps.DimerizedHeisenberg.points_δ_le_one}
  \leanok
  \uses{def:dimerized}
  All three rational doubling certificates of Definition~\ref{def:dimerized} are admissible, every
  \(\beta_0 > 0\), and every \(\delta \le 1\).
\end{lemma}
\begin{proof}
  \leanok
  Verified by exact rational arithmetic from Definition~\ref{def:chain-certificate}.
\end{proof}

\begin{lemma}[Finite-size seed purities for the dimerized spin-\(\tfrac12\) chain]
  \label{lem:dimerized-seeds}
  \lean{CertifiedGaps.DimerizedHeisenberg.seeds}
  \leanok
  \uses{def:dimerized, def:rp-purities}
  At each dimerization in Definition~\ref{def:dimerized}, \(\Ps(4, \beta_0) \ge 1 - s_0\) and
  \(\Pt(8, \beta_0) \ge 1 - t_0\).
\end{lemma}
\begin{proof}
  These finite-size bounds~\eqref{eq:rp-seed-hyp} are certified from complete eigenvalue enclosures
  of the seed rings and enter the formalization as an explicit numerical hypothesis.
\end{proof}

\begin{theorem}[Spectral gap of the dimerized spin-\(\tfrac12\) chain in the rotated frame]
  \label{thm:dimerized-gap}
  \lean{CertifiedGaps.DimerizedHeisenberg.gap}
  \leanok
  \uses{lem:dimerized-valid, lem:dimerized-seeds, thm:rp-gap}
  At each certified dimerization of Definition~\ref{def:dimerized}, every even ring of \(L \ge 4\)
  unit cells in the reflection-positive frame has a unique ground state and a spectral gap greater
  than \(\log r / \beta_0\).
\end{theorem}
\begin{proof}
  \leanok
  Apply Theorem~\ref{thm:rp-gap} with Lemmas~\ref{lem:dimerized-valid} and~\ref{lem:dimerized-seeds}.
\end{proof}

\begin{definition}[Dimerized spin-\(\tfrac12\) Hamiltonian in the laboratory frame]
  \label{def:dimerized-physical}
  \lean{CertifiedGaps.DimerizedHeisenberg.physicalHamiltonian,
    CertifiedGaps.DimerizedHeisenberg.physicalSpectrum, CertifiedGaps.DimerizedHeisenberg.cellRot}
  \leanok
  \uses{def:ring-operators, lem:spin-half-rotation}
  In the laboratory frame, the bond-alternating spin-\(\tfrac12\) chain has Hamiltonian
  \begin{equation}
    \label{eq:dimerized-half-ham}
    H_L = \sum_{j=0}^{L-1} (1+\delta)\, \mathbf S_{2j} \cdot \mathbf S_{2j+1} + (1-\delta)\,
    \mathbf S_{2j+1} \cdot \mathbf S_{2j+2}
  \end{equation}
  written on a ring of strong-bond unit cells, and is related to the reflection-positive frame of
  Definition~\ref{def:dimerized} by applying the local rotation \(1 \otimes e^{i\pi S^y}\) on every
  cell.
\end{definition}

\begin{lemma}[Cellwise unitary equivalence for the dimerized spin-\(\tfrac12\) chain]
  \label{lem:dimerized-physical-eq}
  \lean{CertifiedGaps.DimerizedHeisenberg.physicalHamiltonian_eq,
    CertifiedGaps.DimerizedHeisenberg.physicalSpectrum_eq}
  \leanok
  \uses{def:dimerized-physical, def:dimerized, thm:conj-hamiltonian, thm:conj-spectrum}
  For \(\delta \le 1\) and \(L \ge 2\), the laboratory-frame Hamiltonian~\eqref{eq:dimerized-half-ham}
  is unitarily conjugate to the reflection-positive Hamiltonian of Definition~\ref{def:dimerized},
  so their eigenvalue spectra agree.
\end{lemma}
\begin{proof}
  \leanok
  Rotating the second spin of every cell maps each two-spin exchange term to
  \(-\sum_\alpha P_\alpha \otimes P_\alpha\) by~\eqref{eq:spin-half-rot} in
  Lemma~\ref{lem:spin-half-rotation}; via~\eqref{eq:conj-hamiltonian} in
  Theorem~\ref{thm:conj-hamiltonian}, the inter-cell weak bond becomes
  \((1-\delta) \sum_\alpha (1 \otimes P_\alpha) \otimes (P_\alpha \otimes 1)\), whose nonnegative
  coupling constant equals \(\sqrt{1-\delta}^2 = 1 - \delta\) because \(\delta \le 1\). Applying
  Theorem~\ref{thm:conj-spectrum} completes the proof.
\end{proof}

\begin{theorem}[Uniform spectral gap of the dimerized spin-\(\tfrac12\) Heisenberg ring]
  \label{thm:dimerized-physical-gap}
  \lean{CertifiedGaps.DimerizedHeisenberg.physical_gap}
  \leanok
  \uses{lem:dimerized-physical-eq, thm:dimerized-gap, lem:dimerized-valid}
  At each certified dimerization, every even ring of \(L \ge 4\) unit cells of the dimerized
  spin-\(\tfrac12\) chain~\eqref{eq:dimerized-half-ham} has a unique ground state and spectral gap
  \(> 0.532, 0.172, 0.057\) at \(\delta = 1/2, 3/10, 1/4\).
\end{theorem}
\begin{proof}
  \leanok
  Combine Lemma~\ref{lem:dimerized-physical-eq} with Theorem~\ref{thm:dimerized-gap}.
\end{proof}

\section{Two-leg spin-1/2 Heisenberg ladder}

On a two-leg antiferromagnetic spin-\(\tfrac12\) Heisenberg ladder~\eqref{eq:heisenberg-ladder-ham},
a checkerboard sublattice rotation acts on one spin of every leg bond and one spin of every rung
bond, making the rotated ladder a reflection-positive ring of two-site rungs
(Definition~\ref{def:ladder}, Lemma~\ref{lem:ladder-physical-eq}, and
Theorem~\ref{thm:ladder-physical-gap}).

\begin{definition}[Heisenberg ladder in the reflection-positive frame]
  \label{def:ladder}
  \lean{CertifiedGaps.HeisenbergLadder.rungExchange, CertifiedGaps.HeisenbergLadder.model,
    CertifiedGaps.HeisenbergLadder.Point, CertifiedGaps.HeisenbergLadder.points}
  \leanok
  \uses{def:rpring, def:local-spin-half, def:chain-certificate}
  In the checkerboard-rotated frame, the Heisenberg ladder is a reflection-positive ring
  (Definition~\ref{def:rpring}) of two-site rungs with bond operators \(P_\alpha \otimes 1\) and
  \(1 \otimes P_\alpha\) and rung exchange as the on-site term
  \(o = -J_\perp \sum_\alpha P_\alpha \otimes P_\alpha\). It is certified at rung coupling
  \(J_\perp = 4\) with seed length \(n_0 = 4\) rungs, \(\beta_0 = 2\), and \(r = 9.033\).
\end{definition}

\begin{lemma}[Admissibility of the Heisenberg ladder certificate]
  \label{lem:ladder-valid}
  \lean{CertifiedGaps.HeisenbergLadder.points_valid, CertifiedGaps.HeisenbergLadder.points_β₀_pos}
  \leanok
  \uses{def:ladder}
  The rational doubling certificate of Definition~\ref{def:ladder} is admissible and \(\beta_0 > 0\).
\end{lemma}
\begin{proof}
  \leanok
  Verified by exact rational arithmetic from Definition~\ref{def:chain-certificate}.
\end{proof}

\begin{lemma}[Finite-size seed purities for the Heisenberg ladder]
  \label{lem:ladder-seeds}
  \lean{CertifiedGaps.HeisenbergLadder.seeds}
  \leanok
  \uses{def:ladder, def:rp-purities}
  For the ladder of Definition~\ref{def:ladder}, the seed purities satisfy
  \(\Ps(4, 2) \ge 1 - 0.03509\) and \(\Pt(8, 2) \ge 1 - 0.02463\).
\end{lemma}
\begin{proof}
  These finite-size bounds~\eqref{eq:rp-seed-hyp} are certified from complete eigenvalue enclosures
  of the seed rings and enter the formalization as an explicit numerical hypothesis.
\end{proof}

\begin{theorem}[Spectral gap of the Heisenberg ladder in the rotated frame]
  \label{thm:ladder-gap}
  \lean{CertifiedGaps.HeisenbergLadder.gap}
  \leanok
  \uses{lem:ladder-valid, lem:ladder-seeds, thm:rp-gap}
  Every even ring of \(L \ge 4\) rungs in the reflection-positive frame of Definition~\ref{def:ladder}
  has a unique ground state and \(\Delta_L > \log 9.033 / 2 \approx 1.100\).
\end{theorem}
\begin{proof}
  \leanok
  Apply Theorem~\ref{thm:rp-gap} with Lemmas~\ref{lem:ladder-valid} and~\ref{lem:ladder-seeds}.
\end{proof}

\begin{definition}[Two-leg Heisenberg ladder Hamiltonian in the laboratory frame]
  \label{def:ladder-physical}
  \lean{CertifiedGaps.HeisenbergLadder.legS, CertifiedGaps.HeisenbergLadder.physicalHamiltonian,
    CertifiedGaps.HeisenbergLadder.physicalSpectrum, CertifiedGaps.HeisenbergLadder.rungRot}
  \leanok
  \uses{def:ring-operators, lem:spin-half-rotation}
  In the laboratory frame, the two-leg Heisenberg ladder has Hamiltonian
  \begin{equation}
    \label{eq:heisenberg-ladder-ham}
    H_L = \sum_{\text{legs}} \mathbf S_i \cdot \mathbf S_j + J_\perp \sum_{\text{rungs}}
    \mathbf S_i \cdot \mathbf S_j
  \end{equation}
  and is related to the reflection-positive frame of Definition~\ref{def:ladder} by the
  checkerboard rotation \(e^{i\pi S^y}\) applied to site \((1, j)\) for odd \(j\) and site
  \((2, j)\) for even \(j\).
\end{definition}

\begin{lemma}[Checkerboard unitary equivalence for the Heisenberg ladder]
  \label{lem:ladder-physical-eq}
  \lean{CertifiedGaps.HeisenbergLadder.physicalHamiltonian_eq,
    CertifiedGaps.HeisenbergLadder.physicalSpectrum_eq}
  \leanok
  \uses{def:ladder-physical, def:ladder, thm:conj-hamiltonian, thm:conj-spectrum}
  On a ring with an even number of rungs, the laboratory-frame ladder~\eqref{eq:heisenberg-ladder-ham}
  is unitarily conjugate to the reflection-positive ladder of Definition~\ref{def:ladder}, so their
  eigenvalue spectra agree.
\end{lemma}
\begin{proof}
  \leanok
  On an even ring, the checkerboard rotation acts on one endpoint of every leg bond and one
  endpoint of every rung bond, mapping each exchange interaction to
  \(-\sum_\alpha P_\alpha \otimes P_\alpha\) by~\eqref{eq:spin-half-rot} in
  Lemma~\ref{lem:spin-half-rotation}. Theorems~\ref{thm:conj-hamiltonian}
  and~\ref{thm:conj-spectrum} then complete the proof.
\end{proof}

\begin{theorem}[Uniform spectral gap of the two-leg Heisenberg ladder]
  \label{thm:ladder-physical-gap}
  \lean{CertifiedGaps.HeisenbergLadder.physical_gap}
  \leanok
  \uses{lem:ladder-physical-eq, thm:ladder-gap}
  Every even ring of \(L \ge 4\) rungs of the two-leg Heisenberg ladder
  \eqref{eq:heisenberg-ladder-ham} at \(J_\perp = 4\) has a unique ground state and
  \(\Delta_L > 1.100\).
\end{theorem}
\begin{proof}
  \leanok
  Combine Lemma~\ref{lem:ladder-physical-eq} with Theorem~\ref{thm:ladder-gap}.
\end{proof}

